% Clase 21 — García Olvera Emiliano (actuadores)
\paragraph{Actuadores}
\textbf{Autor:} García Olvera Emiliano

\textbf{En la Pico} (\texttt{actuadores.c}): el motor (GP15) y el buzzer (GP12) usan \texttt{hardware/pwm}, en
\textit{slices} distintos para que el buzzer pueda cambiar de frecuencia sin afectar al motor; los LEDs (GP10 y GP11)
usan \texttt{hardware/gpio}; el apagado lo hace una alarma de \texttt{pico/time} (\texttt{actuadores.c}, líneas 61 y 72),
así ninguna orden bloquea el lazo principal. Son bibliotecas del Pico SDK sobre los registros; no hace falta un driver propio.

\textbf{En la Pi:} las órdenes no usan un dispositivo nuevo, sino el mismo control HID. El puente escribe el reporte de
salida en \texttt{/dev/hidraw}\pend{N} (driver genérico \texttt{hidraw}), y el kernel lo entrega a la Pico como
\texttt{SET\_REPORT} por el \textit{endpoint} de control. En la Pico lo recibe \texttt{tud\_hid\_set\_report\_cb()}
(\texttt{usb\_data\_tx.c}, línea 194).

\textbf{Evidencia:}
\begin{lstlisting}
ls -l /dev/hidraw* /dev/arcade_hid
\end{lstlisting}
\pend{pegar aquí la salida}
